%%%%%%%%%%%%%%%%%%%%%%%%%%%%%%%%%%%%%%%%%%%%%%%%%%%%%%%%%%%%%%%%%%%%%%%%%%%%%%%
%  s01-intro.tex —— 第 1 节 引言（对应原文 pp.1–2）
%%%%%%%%%%%%%%%%%%%%%%%%%%%%%%%%%%%%%%%%%%%%%%%%%%%%%%%%%%%%%%%%%%%%%%%%%%%%%%%

\section{引言}\label{Sec:Introduction}

\begin{keypoints}
\begin{itemize}
  \item \textbf{核心图像}：辛拓扑的出发点虽是经典 Hamilton 动力学，但它自然地出现在代数几何、表示论、
        低维拓扑与弦论中；贯穿这些领域的核心技术是 Floer 上同调。
  \item \textbf{范畴化}：Floer 上同调 $HF^*(L,L')$ 把拉格朗日相交数“范畴化”——
        其 Euler 示性数回到经典的相交数 $[L]\cdot[L']$。
  \item \textbf{结构升级}：$HF^*(L,L)$ 是（有单位元的）环，$HF^*(L,L')$ 是其右模；
        不同拉格朗日之间的 Floer 群由长正合列、拉格朗日对应函子联系在一起，
        并构成 Donaldson 范畴的态射群；在其底下还有一个三角 $A_\infty$-范畴，即 Fukaya 范畴。
  \item \textbf{本文定位}：这是一篇“印象派”式的综述——不追求技术完备，而是解释 Fukaya 范畴是什么、
        从何而来、以及辛拓扑学家想拿它做什么；“计算一个范畴”究竟意味着什么。
\end{itemize}
\end{keypoints}

辛拓扑的源头在于经典 Hamilton 动力学，但它也同样自然地出现在代数几何、表示论、低维拓扑与弦论之中。
在上述所有语境里，辛拓扑所使用的最重要的技术大概就是 \emph{Floer 上同调}：它给 $2n$ 维辛流形
$(X,\omega)$ 中一对定向的 $n$ 维拉格朗日子流形 $L, L'$ 配置一个 $\bZ_2$-分次向量空间
$HF^*(L,L')$，后者把 $L$ 与 $L'$ 的经典同调相交数“范畴化”，其含义是
$\chi(HF^*(L,L')) = (-1)^{n(n+1)/2}\, [L]\cdot[L']$。
% [Tutor Note: 上式中的符号因子 $(-1)^{n(n+1)/2}$ 是（$n$ 维拉格朗日）相交数的标准 Koszul 型符号约定：
%   $n\equiv 0,1 \pmod 4$ 时为正，$n\equiv 2,3 \pmod 4$ 时为负。推导它时要特别注意定向的选取
%   （$L,L'$ 的定向与 $X$ 的定向如何相容），此处极易出现整体符号错误；
%   核对时请从 $T_pX \cong T_pL\oplus T_pL'$ 的定向次序入手。]

在这门学科的早期，Floer 上同调所扮演的角色，类似于上世纪之初奇异同调群在代数拓扑中所扮演的角色：
这些群的秩为几何或动力系统来源的问题（例如周期点的个数）给出下界。但是，正如它的经典代数拓扑同胞一样，
随着 Floer 上同调获得了越来越多的结构，它的重要性也与日俱增。对 $(L,L)$ 而言，Floer 上同调群构成一个
有单位元的环，而 $HF^*(L,L')$ 是这个环上的右模；不同拉格朗日所对应的 Floer 群常常被长正合列串在一起；
这些群在由拉格朗日对应所定义的辛流形态射下具有函子性；并且它们本身就是一个以拉格朗日子流形为对象、
与 $X$ 相联系的范围——Donaldson 范畴——中的态射群。

% [Tutor Note: 这里出现了三个层次，务必分清：链层次的 $CF^*(L,L')$；上同调层次的 $HF^*(L,L')$
%   （即 Donaldson 范畴的态射）；以及下面要说的带高阶运算 $\mu^k$ 的 $A_\infty$-范畴（Fukaya 范畴）。
%   初学者最常犯的错误是把 “Fukaya 范畴” 与 “Donaldson 范畴” 混为一谈：前者是后者的
%   $A_\infty$-增强，上同调取 $H^0$ 后才回到后者。]

近来，人们的注意力集中于更精细的代数结构，以及 Donaldson 范畴底下那个三角 $A_\infty$-范畴
$\scrF(X)$，即所谓的 \emph{Fukaya 范畴}。$A_\infty$-结构在数学的其他地方也屡见不鲜；而且——
很大程度上是在 Kontsevich 那惊人的\emph{同调镜像对称}猜想的推动之下——人们现在认识到，
其他领域的研究者将会“对着镜子观看，模糊不清”地瞥见辛拓扑，其原因恰恰在于：那些领域中
自然受到关注的范畴，可以被解释为辛流形的 Fukaya 范畴（而且常常是些乍看之下在那“姐妹理论”中
根本无迹可寻的辛流形）。镜像对称的表述使代数几何学家得以看清 Fukaya 范畴与他们所研究的
层的导出范畴之间的关联；除此之外，它与对 Langlands 纲领感兴趣的几何表示论学家、
研究箭图变异的丛代数学家、以及在规范理论或纽结量子不变量方面工作的低维拓扑学家之间，
也都有联系，如此等等。说来有些反常：对某些人而言，Fukaya 范畴倒可能是他们与辛拓扑的
第一个接触点。

% [Tutor Note: “through a glass darkly”（对着镜子观看，模糊不清）出自《哥林多前书》13:12，
%   作者借此形容镜像对称两侧的对应关系是间接的、需要通过层层翻译才能看清的。]

对于有心把 Fukaya 范畴纳入日常工具箱的研究生，目前既有实质性的系统论述 \cite{FO3, FCPLT}，
也有较为平易的导引 \cite{Auroux}。本文这篇更偏“印象派”的综述，或许仍不无裨益。我们希望它
\emph{能够}促使那些使用全纯曲线理论的人，去认真对待采取范畴视角所需要付出的额外努力：
我们试图说明，代数复杂性的增加如何扩展了理论的适用范围，以及“计算”一个范畴究竟意味着什么。
更宽泛地说，我们希望让 Fukaya 范畴对那些（也许身处相邻领域的）希望了解辛拓扑关心什么、
三角范畴如何在该领域中产生、以及我们可以指望用它们做些什么的人，变得更容易接近——
至少是更熟悉。

\noindent\textbf{致歉.} Floer 理论充满了各种微妙之处，我们或者对它们置之不理\footnote{%
\emph{作品的删节者伤害知识}——列奥纳多·达·芬奇，约 1510 年。}，
或者把它们指认出来然后一掠而过，只求不遮蔽更宏阔的图景\footnote{%
感谢匿名审稿人，使笔者“滑冰的姿势”得到了若干改进。}。
这是一个迅速演进的领域，本文所给出的快照远非全面。

\vspace{0.2cm}
\noindent\textbf{致谢.} 笔者对 Fukaya 范畴的看法，主要是在与 Mohammed Abouzaid、Denis Auroux
和 Paul Seidel 的交谈与合作中形成的。能够追随他们的工作学习这一学科，是笔者的荣幸。

% [Tutor Note: 原文标题 “Apologies.” 与 “Acknowledgements” 是作者的自谦笔法。
%   “skating over”（一掠而过/滑过去）与审稿人的 “improvements to my skating style”
%   构成一个双关：既指滑冰，也指在技术细节上“滑过不提”。译文保留了引号以示意。]
